\section{Escape effects and scope-minimal normal form}
\label{sec:effects}

\subsection{Static effect information}

A proof graph node that introduces a certificate is annotated with one of the
following effects.

\begin{center}
\begin{tabularx}{0.94\textwidth}{@{}lXl@{}}
\toprule
Effect & Permitted public dependence & Required evidence \\
\midrule
\(\mathsf{sealed}\) & none & producer or residual; local extractor \\
\(\mathsf{rel}(O,\approx)\) & existential observation relation modulo \(\approx\) & producer or residual; coherence proof \\
\(\mathsf{named}(p)\) & concrete PR term \(p(q)\) & direct validity proof \\
\(\mathsf{named}(\mu,B)\) & least accepted code below \(B(q)\) & reflected verifier and bounded existence \\
\bottomrule
\end{tabularx}
\end{center}

The first two modes do not move a certificate eigenvariable to the outer
scope.  The named modes replace it with an outer term already defined without
that eigenvariable.

A host-side checker also records whether a missing producer has a named
residual.  A node with no producer and no residual is rejected; otherwise the
residual becomes an explicit context formula in the compiled theorem.

\subsection{Scope regions}

For an acyclic dependency graph, the \emph{consumer region} of a certificate
node is the subgraph consisting of all nodes whose local derivations depend on
that certificate before it is sealed or converted to a justified public
output.  Opening the witness outside this region is harmless but bloats the
proof and may force unrelated code bounds.  Opening it inside the region would
leave some consumer without the witness.

A scope-minimal plan opens the certificate immediately before the least proof
region containing all its consumers and closes it immediately after that
region's result has become certificate-independent.

\begin{theorem}[Scope-minimal normal form]
\label{thm:scope-normal-form}
Let \(G\) be a finite acyclic proof graph in which every certificate node has a
valid escape effect and every local edge is backed by an object-iPRA derivation.
Then \(G\) can be transformed effectively into a nested proof plan with the
following properties.

\begin{enumerate}[label=(\roman*)]
\item Every sealed certificate is introduced by one existential elimination
  surrounding exactly its live consumer region.
\item Every relational certificate is hidden inside its public output relation;
  no representative occurs free outside that relation's eliminator.
\item Every named certificate has been replaced by a direct producer term or a
  least bounded selector before it appears in an outer formula.
\item Unreachable nodes and their producer residuals are absent.
\item All remaining unresolved producer formulas occur in the outer context.
\end{enumerate}
\end{theorem}

\begin{proof}
Choose a reverse topological order of the graph.  Remove every node not
reachable from the final goal.  For a live certificate \(c\), collect all edges
whose derivations require its validity formula.  Replace the finite collection
of downstream conclusions by one continuation deriving their conjunction or
the final joined result under the single local assumption \(V(c)\).  This is
possible because weakening permits unrelated local assumptions and ordinary
conjunction introduction shares the witness inside one scope.

Now move the existential elimination inward across any inference whose other
premises and conclusion do not contain the certificate.  The standard
commuting conversion is valid because the eigenvariable is fresh there.  Each
move strictly decreases the number of graph edges crossed by the scope, so the
finite process terminates at the least consumer region.  Repeat in reverse
nesting order for dependent certificates.

For a relational effect, replace the certificate-dependent observer term in
the outer interface by the existential relation \(\Out\); the proofs of
\cref{thm:observed-relation} justify all public comparisons.  For a named
effect, insert the supplied producer or \cref{thm:least-code}, yielding a term
whose free variables are outer parameters only.  Finally move unresolved
producer assumptions outward by implication/context introduction.  Every
transformation is a composition of weakening, cut, conjunction rules,
existential rules, and substitution, so the resulting plan has the same
object-theory conclusion.
\end{proof}

\subsection{No-escape safety}

\begin{corollary}[No dangling certificate]
A well-effected plan cannot compile to a public formula containing a proof-local
certificate variable.
\end{corollary}

\begin{proof}
In the sealed and relational cases the certificate is bound by an existential
inside the plan.  In the named cases it is replaced by an outer term.  The
indexed type of existential elimination requires the final formula to live at
the outer arity, so no unbound local variable can occur.
\end{proof}

In Lean this property does not need a separate free-variable checker.  It is a
consequence of the indices of \code{Fm}, contexts, and \code{Prf.exE}.  The
host effect checker adds a mathematical policy: it rejects an exact public
representative when the formal scope could only support a relation.

\subsection{Master certificates as an optional backend}

If every internal witness has an explicit finite bound, the nested plan can be
packed into one mixed-radix master code.  Its checker is the conjunction of all
local checks, and bounded search selects the least accepted master code.
Conversely, a master code can be opened once and its fields used in nested
scopes.  Thus bounded scoped and packed presentations are extensionally
intertranslatable.

The scoped representation is strictly more general as an interface: it remains
meaningful when a proof-local existential witness has no useful numerical
bound.  It also avoids changing every radix when one local certificate format
changes.
